\documentclass[12pt,a4paper]{article}

\usepackage[utf8]{inputenc}
\usepackage[T1]{fontenc}
\usepackage[portuguese]{babel}

\usepackage[colorlinks=true, urlcolor=blue, linkcolor=black]{hyperref}
\usepackage{geometry}
\usepackage{textcomp}

\geometry{
    a4paper,
    left=3cm,
    right=3cm,
    top=3cm,
    bottom=3cm,
}

\begin{document}

\begin{center}
    {\Large \textbf{Bibliografia - LoRa (Long Range) \& LoRaWAN}} \\[0.6cm]
    {\large \textbf{Danilo Davi Gomes Fróes}} \\[0.4cm]
    {\large \textbf{Grupo 18}}
\end{center}

\vspace{1cm}

\noindent
SEMTECH CORPORATION. AN1200.22: LoRa\texttrademark\ Modulation Basics. Application Note, 2015. Disponível em: \href{http://frugalprototype.com/wp-content/uploads/2016/08/an1200.22.pdf}{frugalprototype.com/wp-content/uploads/2016/08/an1200.22.pdf}

\vspace{0.6cm}
\noindent
LORA ALLIANCE. LoRaWAN\textregistered\ Specification v1.1. 2020. Disponível em: \\ \href{https://resources.lora-alliance.org/technical-specifications/lorawan-specification-v1-1}{resources.lora-alliance.org/technical-specifications/lorawan-specification-v1-1}

\vspace{0.6cm}
\noindent
ADELANTADO, Ferran, et al. Understanding the Limits of LoRaWAN. IEEE Communications Magazine. 55. 10.1109/MCOM.2017.1600613. 2017.

\vspace{0.6cm}
\noindent
BOR, Martin C., et al. Do LoRa Low-Power Wide-Area Networks Scale?. In: Proceedings of the 19th ACM International Conference on Modeling, Analysis and Simulation of Wireless and Mobile Systems (MSWiM), 2016.

\vspace{0.6cm}
\noindent
THE THINGS NETWORK (TTN). LoRaWAN Documentation. Disponível em: \\ \href{https://www.thethingsnetwork.org/docs/lorawan/}{www.thethingsnetwork.org/docs/lorawan/}

\vspace{0.6cm}
\noindent
BERTOLETI, Pedro. Conectividade LoRaWAN: Fundamentos e Prática. São Paulo: Editora NCB, 2023.

\end{document}